\chapter{计算机组成基础}

\begin{introduction}
  \item 计算机由哪几个部分组成
  \item 程序到底是什么
  \item 计算机为什么只认识 0 和 1
  \item 内存、地址和变量
  \item 从写代码到程序运行
\end{introduction}

\chaptermotto{「计算机很笨，但它很听话。」}

同学们，先别急着写代码。我们先用一节课的时间，认识一下计算机这个“大玩具”。

你可能会觉得计算机很聪明：它会算题、会放视频、会下棋。其实它一点都不聪明。它只会做很简单的事情：\textbf{算一算、比一比、把东西搬来搬去}。它之所以看起来聪明，是因为人把做事情的步骤一条一条写给了它。

这一章不写程序，只讲五件事：计算机由什么组成、程序是什么、为什么它只认识 0 和 1、内存是什么、以及你写的代码是怎么变成能运行的程序的。

\section{计算机能做什么}

我们用计算机做事情，通常都是这样一个过程：

\begin{center}
\begin{tikzpicture}[x=1cm,y=1cm,>=Stealth,node distance=12mm]
  \node[box] (in) {输入};
  \node[box,right=of in] (do) {处理};
  \node[box,right=of do] (out) {输出};
  \draw[flow] (in) -- (do);
  \draw[flow] (do) -- (out);
\end{tikzpicture}
\end{center}

举几个例子：

\begin{itemize}
  \item 用计算器算 $3+5$：你按的“3、+、5”是\textbf{输入}；计算器在内部做加法是\textbf{处理}；屏幕上出现“8”是\textbf{输出}。
  \item 用手机拍照：镜头看到画面是输入；手机把画面变成照片是处理；屏幕上显示照片是输出。
  \item 让小车前进：传感器测到距离是输入；单片机判断“前面没有障碍”是处理；电机转动是输出。
\end{itemize}

\textbf{输入 → 处理 → 输出}，这就是计算机做事的三个步骤。我们写程序，主要就是在告诉计算机“怎么处理”。

\section{计算机的四样东西}

一台计算机看起来零件很多，但真正重要的只有四样：

\begin{figure}[htbp]
  \centering
  \begin{tikzpicture}[x=1cm,y=1cm,>=Stealth]
    \node[cpubox,minimum width=3cm,minimum height=1.4cm] (cpu) at (0,1.8) {CPU\\（大脑）};
    \node[box,minimum width=2.8cm] (mem) at (-4,-0.8) {内存\\（草稿纸）};
    \node[box,minimum width=2.8cm] (disk) at (4,-0.8) {硬盘\\（书包）};
    \node[box,minimum width=2.8cm] (io) at (0,-2.8) {输入 / 输出\\（眼睛、嘴巴）};

    \draw[<->,>=Stealth,draw=structurecolor,line width=0.8pt] (cpu.west) -| (mem.north);
    \draw[<->,>=Stealth,draw=structurecolor,line width=0.8pt] (cpu.east) -| (disk.north);
    \draw[<->,>=Stealth,draw=structurecolor,line width=0.8pt] (cpu.south) -- (io.north);
    \draw[<->,>=Stealth,draw=structurecolor,line width=0.8pt] (mem.east) -- (io.west);
    \draw[<->,>=Stealth,draw=structurecolor,line width=0.8pt] (disk.west) -- (io.east);
  \end{tikzpicture}
  \caption{计算机的四样核心部件：CPU、内存、硬盘、输入输出}
  \label{fig:four-parts}
\end{figure}

\subsection{CPU：计算机的大脑}

CPU 是计算机的“大脑”，全名叫\textbf{中央处理器}。

它负责两件事：

\begin{enumerate}
  \item \textbf{算}：做加减乘除、比较大小、判断真假；
  \item \textbf{指挥}：告诉其他部件“现在该做什么”。
\end{enumerate}

CPU 算得很快，一秒钟能做几十亿次运算。但它只会做很简单的运算，复杂的事情要拆成很多小步骤。

\subsection{内存：计算机的草稿纸}

内存是计算机的“草稿纸”，用来临时存放正在使用的数据。

它的特点是：

\begin{itemize}
  \item 速度比较快；
  \item \textbf{一断电，里面的东西就全没了}。
\end{itemize}

所以内存只适合放“正在用”的东西。比如你正在写的文档、正在玩的游戏，数据都在内存里。突然断电，没保存的内容就丢了。

\subsection{硬盘：计算机的书包}

硬盘是计算机的“书包”，用来长期存放东西。

它的特点是：

\begin{itemize}
  \item 速度比内存慢；
  \item \textbf{断电以后东西还在}。
\end{itemize}

你的照片、作业、程序文件，平时都放在硬盘里。用的时候，计算机会把它们从硬盘搬到内存里。

\begin{tipbox}[手机也是一样的]
手机的“8 GB + 256 GB”里，8 GB 通常指运行内存（相当于电脑的内存），256 GB 是存储空间（相当于硬盘）。内存越大，能同时开的程序越多；存储越大，能存的东西越多。
\end{tipbox}

\subsection{输入和输出：眼睛、耳朵和嘴巴}

输入设备让计算机“知道外面发生了什么”，例如键盘、鼠标、麦克风、摄像头、温度传感器。

输出设备让计算机“把结果告诉外面”，例如屏幕、音箱、打印机、电机、指示灯。

\begin{center}
\begin{tabular}{@{}lll@{}}
\toprule
部件 & 像人的什么 & 作用 \\
\midrule
CPU & 大脑 & 计算、指挥 \\
内存 & 草稿纸 & 临时存放正在用的数据 \\
硬盘 & 书包 & 长期存放数据 \\
输入设备 & 眼睛、耳朵 & 把信息送进计算机 \\
输出设备 & 嘴巴、手 & 把结果送出去 \\
\bottomrule
\end{tabular}
\end{center}

\section{程序是什么}

\subsection{程序就是一张步骤清单}

程序不是魔法，它只是一张\textbf{步骤清单}。就像泡面的做法：

\begin{enumerate}
  \item 烧一壶水；
  \item 把面饼放进碗里；
  \item 倒入开水；
  \item 等三分钟；
  \item 放入调料，搅拌。
\end{enumerate}

计算机做事情也是这样：你给它一张步骤清单，它就从第一条开始，一条一条做下去。

\subsection{计算机不会猜}

人听别人说话，可以猜意思。计算机不会猜。你写“把水烧开”，它不知道多少度算开、用什么烧、要等多久。你必须写清楚：

\begin{itemize}
  \item 用哪个锅？
  \item 加热到多少度？
  \item 到了温度以后做什么？
\end{itemize}

写程序也是一样：\textbf{每一步都要明确，不能有歧义}。这就是为什么程序设计要学很多规则。

\begin{tipbox}[把计算机当成一个非常听话、但非常死板的同学]
你说“往前走”，他就一直往前走，撞墙也不停，因为你没说“遇到墙就停”。写程序时，要把所有情况都想到。
\end{tipbox}

\section{计算机为什么只认识 0 和 1}

\subsection{电灯只有开和关}

计算机里面是电子电路。电路最简单的两种状态是：

\begin{itemize}
  \item 通电；
  \item 断电。
\end{itemize}

就像电灯：要么亮，要么灭。我们用数字表示这两种状态：

\[
\text{通电} = 1, \qquad \text{断电} = 0.
\]

只用两种状态，电路最简单、最不容易出错。所以计算机选择用 0 和 1 来表示所有信息。

\subsection{位和字节}

\begin{itemize}
  \item 一个 0 或一个 1，叫一位（bit，比特）；
  \item 8 位合在一起，叫一个字节（byte）。
\end{itemize}

\begin{center}
\begin{tabular}{@{}cccccccc@{}}
\toprule
1 & 0 & 1 & 1 & 0 & 0 & 1 & 0 \\
\bottomrule
\end{tabular}
\end{center}

这是一个字节，一共 8 位。一个字节可以组合出 $2^8 = 256$ 种不同的样子（从 \texttt{00000000} 到 \texttt{11111111}）。

\subsection{用 0 和 1 数数}

我们平时用十进制：0、1、2、3、4、5、6、7、8、9，到 9 再往上就进位，变成 10。

计算机用二进制：只有 0 和 1，到 1 再往上就进位：

\begin{center}
\begin{tabular}{@{}ll@{}}
\toprule
二进制 & 十进制 \\
\midrule
\texttt{0} & 0 \\
\texttt{1} & 1 \\
\texttt{10} & 2 \\
\texttt{11} & 3 \\
\texttt{100} & 4 \\
\texttt{101} & 5 \\
\texttt{110} & 6 \\
\texttt{111} & 7 \\
\texttt{1000} & 8 \\
\bottomrule
\end{tabular}
\end{center}

规律是：\textbf{满二进一}。和我们数数“满十进一”是一样的，只是计算机只有两个数字。

\begin{example}
二进制 \texttt{1011} 表示十进制多少？
\end{example}
\begin{solution}
从右往左，每一位的“份量”分别是 1、2、4、8：

\[
\texttt{1011} = 8 + 0 + 2 + 1 = 11.
\]

所以二进制 \texttt{1011} 就是十进制的 11。
\end{solution}

\subsection{字符怎么表示}

计算机里没有“字母”，只有数字。那怎么存字母呢？

人们约定了一张“密码表”，给每个字符编一个数字。最常用的是 \textbf{ASCII 码}。例如：

\begin{center}
\begin{tabular}{@{}lll@{}}
\toprule
字符 & ASCII 码 & 说明 \\
\midrule
\texttt{'0'} & 48 & 数字字符 0 \\
\texttt{'A'} & 65 & 大写字母 A \\
\texttt{'a'} & 97 & 小写字母 a \\
空格 & 32 & 空格 \\
\texttt{'\textbackslash n'} & 10 & 换行 \\
\bottomrule
\end{tabular}
\end{center}

所以计算机里存的不是字母 \texttt{A}，而是数字 65。程序输出时，再按照密码表把 65 变回字母 \texttt{A}。

汉字也有编码表，例如 GBK、UTF-8。一个汉字通常要占 2 个或 3 个字节。这部分以后用到再了解。

\begin{tipbox}[负数和小数怎么办？]
负数和小数在计算机里也有专门的表示方法，但比正数复杂。初学阶段先不用管，只要知道“计算机有办法表示”就可以了。等你把基本程序写熟，再回头看也不迟。
\end{tipbox}

\section{内存像一排小盒子}

\subsection{地址就是盒子编号}

内存可以想象成一排小盒子，每个盒子装 1 个字节。每个盒子都有一个编号，这个编号叫\textbf{地址}。

\begin{figure}[htbp]
  \centering
  \begin{tikzpicture}[x=1cm,y=1cm]
    \foreach \i/\addr/\v in {0/1000/?,1/1001/?,2/1002/?,3/1003/?,4/1004/?} {
      \node[font=\ttfamily\scriptsize,color=black!55,anchor=east] at (-0.2,-\i*0.9) {\addr};
      \node[membox,minimum width=2.2cm] at (1.6,-\i*0.9) {\v};
    }
    \draw[decorate,decoration={brace,amplitude=4pt},structurecolor!70]
      (3.0,0.3) -- (3.0,-2.4);
    \node[anchor=west,font=\small] at (3.2,-1.05) {从 1000 号盒子开始};
  \end{tikzpicture}
  \caption{内存像一排有编号的小盒子，每个盒子有地址}
  \label{fig:memory-boxes}
\end{figure}

地址就像酒店的房间号：房间号不是房间本身，但拿着房间号就能找到房间。

\subsection{变量就是贴了标签的盒子}

写程序时，我们不需要每次都记“1000”这样的地址。我们给盒子贴一个标签，起一个名字，例如：

\begin{codebox}{variable-box.c —— 变量就是一个盒子}
int a = 10;
\end{codebox}

这句话的意思是：

\begin{itemize}
  \item 准备一个能装整数的盒子；
  \item 在盒子外面贴上标签 \texttt{a}；
  \item 往盒子里放 10。
\end{itemize}

以后写 \texttt{a}，就是在说“那个贴了 \texttt{a} 标签的盒子”。盒子里的东西可以换，但标签不变。

\begin{tipbox}[变量三件事]
每个变量都有：\textbf{名字}（标签）、\textbf{类型}（盒子能装什么）、\textbf{值}（盒子里现在装什么）。地址是盒子的编号，通常由计算机管理，不需要你手动记。
\end{tipbox}

\section{从写代码到程序运行}

\subsection{人话和机器话}

你写的 C 代码，长这样：

\begin{codebox}{hello.c —— 给人看的代码}
#include <stdio.h>

int main(void)
{
    printf("Hello!\n");
    return 0;
}
\end{codebox}

但 CPU 看不懂英文，它只认识 0 和 1 组成的机器指令。所以中间需要一个“翻译官”。

\subsection{编译器就是翻译官}

把 C 代码翻译成机器能懂的程序，这个工具叫\textbf{编译器}。

\begin{figure}[htbp]
  \centering
  \begin{tikzpicture}[x=1cm,y=1cm,>=Stealth,node distance=13mm]
    \node[box,minimum width=2.8cm] (src) {你写的代码\\\texttt{hello.c}};
    \node[box,minimum width=2.8cm,right=of src] (cc) {编译器\\（翻译官）};
    \node[box,minimum width=2.8cm,right=of cc] (exe) {可执行程序\\\texttt{hello.exe}};
    \node[box,minimum width=2.8cm,right=of exe] (run) {双击运行\\屏幕上出字};
    \draw[flow] (src) -- (cc);
    \draw[flow] (cc) -- (exe);
    \draw[flow] (exe) -- (run);
  \end{tikzpicture}
  \caption{写代码 → 编译 → 运行}
  \label{fig:compile-simple}
\end{figure}

过程是这样的：

\begin{enumerate}
  \item 你在电脑上写一个 \texttt{hello.c} 文件；
  \item 用编译器把它翻译成 \texttt{hello.exe}；
  \item 双击 \texttt{hello.exe}，计算机就按里面的指令做事。
\end{enumerate}

\subsection{两种错误}

写程序会遇到两种错误：

\begin{itemize}
  \item \textbf{编译错误}：翻译官看不懂你写的话。比如少了一个分号、括号不配对。这时程序根本生成不了。
  \item \textbf{运行错误}：程序能生成，但运行的时候做错了事。比如让计算机除以 0，或者读到了不该读的地方。
\end{itemize}

\begin{tipbox}[报错不可怕]
编译器报错时，会告诉你“第几行出了问题”。先看第一行错误，通常改掉它，后面的错误就消失了。
\end{tipbox}

\section{C 语言为什么适合电控}

你以后要接触电机、传感器、单片机、机器人。这些设备有三个特点：

\begin{itemize}
  \item 要\textbf{快}：电机控制一秒钟要算很多次；
  \item 要\textbf{省}：单片机的内存很小；
  \item 要\textbf{直接控制硬件}：读写寄存器、引脚、串口。
\end{itemize}

C 语言正好适合：

\begin{enumerate}
  \item 它很\textbf{接近硬件}，能直接操作内存和寄存器；
  \item 它\textbf{很快}，没有太多额外的开销；
  \item 它\textbf{通用}，从单片机到电脑都能用；
  \item 学好 C，再学 C++、Python、数据结构都会更容易。
\end{enumerate}

\section{本章小结}

\begin{center}
\begin{tabular}{@{}p{0.28\textwidth}p{0.6\textwidth}@{}}
\toprule
问题 & 一句话答案 \\
\midrule
计算机做什么 & 输入 → 处理 → 输出 \\
有哪四样核心部件 & CPU、内存、硬盘、输入输出 \\
CPU 是什么 & 大脑，负责计算和指挥 \\
内存和硬盘的区别 & 内存快但断电就丢；硬盘慢但能长期保存 \\
程序是什么 & 给计算机的步骤清单，必须一步一步写清楚 \\
计算机为什么用二进制 & 电路只有通/断两种稳定状态 \\
位和字节 & 1 位是 0 或 1；8 位是 1 字节 \\
内存是什么 & 一排有编号的小盒子，编号叫地址 \\
变量是什么 & 贴了名字标签的盒子 \\
代码怎么运行 & 写代码 → 编译器翻译 → 可执行程序 → 运行 \\
\bottomrule
\end{tabular}
\end{center}

\begin{tipbox}[下一章预告]
认识完计算机，下一章我们就写第一个真正的 C 程序：让计算机在屏幕上说话，再让它听我们输入数字、做计算。
\end{tipbox}

\begin{problemset}
  \item 用自己的话说一说：计算机做一件事要经过哪三个步骤？
  \item CPU、内存、硬盘各像什么？它们分别负责什么？
  \item 内存和硬盘最大的区别是什么？为什么手机内存越大越贵？
  \item 程序为什么必须一步一步写清楚？举一个“没说清楚”的例子。
  \item 二进制 \texttt{1010}、\texttt{1100}、\texttt{11111} 分别等于十进制多少？
  \item 1 个字节等于多少位？一个字节最多能表示多少种不同的组合？
  \item 字符 \texttt{'A'} 在计算机里存的是什么？为什么不是直接存字母？
  \item 把“写代码 → 运行”的过程画出来，并说明编译器做了什么。
  \item 编译错误和运行错误有什么区别？各举一个例子。
  \item 思考题：为什么计算机不直接用十进制，而是用二进制？
\end{problemset}
